\documentclass[11pt]{amsart}
\usepackage[T1]{fontenc}
\usepackage{lmodern,microtype,amsmath,amssymb,amsthm,mathtools,booktabs}
\usepackage[margin=1.1in]{geometry}
\usepackage[hidelinks]{hyperref}
\usepackage{xurl}

\newtheorem{theorem}{Theorem}[section]
\newtheorem{lemma}[theorem]{Lemma}
\newtheorem{corollary}[theorem]{Corollary}
\theoremstyle{remark}
\newtheorem{remark}[theorem]{Remark}
\numberwithin{equation}{section}

\newcommand{\E}{\mathbb{E}}
\renewcommand{\P}{\mathbb{P}}
\newcommand{\one}{\mathbf{1}}
\DeclareMathOperator{\Var}{Var}
\DeclareMathOperator{\tr}{tr}

\hypersetup{pdftitle={The complete graph maximizes the variance of the antivoter model},pdfauthor={Yitzchak Shmalo}}

\title[The antivoter variance]{The complete graph maximizes the variance of the antivoter model}
\author{Yitzchak Shmalo}
\address{Einstein Institute of Mathematics, The Hebrew University of Jerusalem, Jerusalem, Israel}
\email{yitzchak.shmalo@gmail.com}
\date{8 November 2026}

\begin{document}

\begin{abstract}
In the antivoter model on a finite graph each vertex, at rate one, adopts the opinion opposite to that of a uniformly
chosen neighbour. Let $S$ be the sum of the $\pm1$ opinions in stationarity. Aldous and Fill asked whether, among
regular graphs on $n$ vertices, $\Var S$ is largest for the complete graph, where it equals $n(n-2)/(2n-3)$. We prove
that it is: for every connected non-bipartite regular graph on $n\ge3$ vertices and every stationary distribution,
$\Var S\le n(n-2)/(2n-3)$, with equality only for $K_n$. The proof writes the pair correlations as the solution of a
Sylvester equation and reduces the question to a lower bound on the diagonal of $(I+P)^{-1}$, where $P$ is the
transition matrix of simple random walk. For vertex-transitive graphs the same argument gives
$\Var S=n^2/\bigl(2\sum_i(1+\lambda_i)^{-1}\bigr)$, the sum running over the eigenvalues of $P$.
\end{abstract}

\maketitle

\section{Introduction}

Let $G=(V,E)$ be a finite connected $d$-regular graph on $n\ge3$ vertices. The antivoter model on $G$ is the
continuous-time Markov chain on $\{-1,1\}^V$ in which each vertex $v$, at the times of a rate-one Poisson clock,
chooses a neighbour $u$ uniformly at random and sets $\eta_v\coloneqq-\eta_u$. On finite graphs it was studied by
Donnelly and Welsh~\cite{DW}, who showed that when $G$ is neither bipartite nor a cycle the chain is
irreducible on the $2^n-2$ non-constant configurations; on a bipartite graph it is absorbed in one of the two proper
$2$-colourings. Write $S=\sum_v\eta_v$ for the opinion margin under a stationary distribution. The antivoter model is a
standard test case for normal approximation by Stein's method~\cite{RR,Ross}, where $\Var S$ enters as the
normalising constant.

Section~14.4 of Aldous and Fill~\cite{AF} develops the variance theory: $\E S=0$, $\Var S<n/2$ for every such graph
(their Proposition~14.22), $2\kappa/d\le\Var S$ with $\kappa$ a cut parameter (Corollary~14.21), and for the complete
graph the correlation of two distinct vertices is $-1/(2n-3)$, so that
\begin{equation}\label{eq:Kn}
\Var S=\frac{n(n-2)}{2n-3}\qquad (G=K_n).
\end{equation}
They then ask for a proof that $\Var S$ is maximized by the complete graph, and remark that this might follow from a
sharpening of the inequality behind $\Var S<n/2$. We answer the question.

\begin{theorem}\label{thm:main}
Let $G$ be a connected non-bipartite $d$-regular graph on $n\ge3$ vertices, and let $\eta$ have any stationary
distribution of the antivoter model on $G$. Then
\[
\Var S\le\frac{n(n-2)}{2n-3},
\]
with equality if and only if $G=K_n$.
\end{theorem}

The hypothesis that $G$ is not a cycle, needed for irreducibility, plays no role: the bound holds for every stationary
distribution, and the odd cycles are covered. Bipartite graphs are excluded because there the stationary distributions are
mixtures of the two proper colourings, on which $S=\pm(|A|-|B|)$ for the colour classes $A,B$, so $\Var S$ can be of
order $n^2$.

The argument also yields an exact formula for vertex-transitive graphs. Let $P=A_G/d$ be the transition matrix of
simple random walk, with eigenvalues $1=\lambda_1>\lambda_2\ge\dots\ge\lambda_n>-1$.

\begin{corollary}\label{cor:vt}
If $G$ is in addition vertex-transitive, then
\[
\Var S=\frac{n^2}{2\sum_{i=1}^n(1+\lambda_i)^{-1}}.
\]
\end{corollary}

For $K_n$ the eigenvalues are $1$ and $-1/(n-1)$ with multiplicity $n-1$, and the corollary returns~\eqref{eq:Kn}. For
the octahedron $K_{2,2,2}$ ($\lambda=1,0,0,0,-\tfrac12,-\tfrac12$) it gives $36/15=2.4$, and for the prism $C_3\times
K_2$ ($\lambda=1,\tfrac13,0,0,-\tfrac23,-\tfrac23$) it gives $36/18.5\approx1.9459$; both agree with the exact solution
of the correlation equations in Table~\ref{tab:small}.

\section{The correlation equations}

Fix a stationary distribution and put $c(v,w)=\E[\eta_v\eta_w]$, so $c(v,v)=1$, and let $C=(c(v,w))_{v,w\in V}$. Let
$Q=I+P$. Since $G$ is connected and not bipartite, the eigenvalues of $P$ lie in $(-1,1]$, so $Q$ is symmetric and
positive definite with $Q\one=2\one$.

\begin{lemma}\label{lem:eq}
$\E\eta_v=0$ for every $v$, and there are numbers $\delta_v\ge0$ such that
\begin{equation}\label{eq:syl}
QC+CQ=D,\qquad D=\operatorname{diag}(\delta_v),\qquad \delta_v=2\Bigl(1+\frac1d\sum_{u\sim v}c(u,v)\Bigr).
\end{equation}
Moreover $\Var S=\tfrac14\sum_v\delta_v$.
\end{lemma}

\begin{proof}
Under stationarity the generator applied to any function has mean zero. For $f(\eta)=\eta_v$ only the clock at $v$
changes $f$, and it replaces $\eta_v$ by $-\eta_u$ with $u$ a uniform neighbour, so $0=-\E\eta_v-(Pm)_v$ with
$m_v=\E\eta_v$. Thus $Qm=0$ and $m=0$. For $f(\eta)=\eta_v\eta_w$ with $v\ne w$, the clocks at $v$ and at $w$ each change
$f$, and
\[
0=-2c(v,w)-\sum_{u}P(v,u)c(u,w)-\sum_{u}P(w,u)c(v,u),
\]
which is the $(v,w)$ entry of $QC+CQ=0$ off the diagonal (the term $u=w$ in the first sum is $c(w,w)=1$, as it should
be, since $\eta_w\eta_w=1$). Hence $QC+CQ$ is diagonal; its $(v,v)$ entry is $2(QC)_{vv}=2(1+(PC)_{vv})$, which is the
$\delta_v$ of~\eqref{eq:syl}, and $\delta_v\ge0$ because $|c|\le1$. Finally, $\E S=0$ by the first part, and
multiplying~\eqref{eq:syl} by $\one$ on both sides gives $4\,\one^{\top}C\one=\sum_v\delta_v$, since $Q\one=2\one$ and
$Q$ is symmetric; and $\one^{\top}C\one=\E S^2=\Var S$.
\end{proof}

In terms of correlations, Lemma~\ref{lem:eq} says $\Var S=\frac n2(1+\bar c)$, where $\bar c$ is the average of
$c(u,v)$ over the $nd/2$ edges; this is Lemma~14.20 of~\cite{AF} in another form, and Corollary~14.24 there when $G$ is
edge-transitive.

\begin{lemma}\label{lem:rep}
With $m_u=(Q^{-1})_{uu}$,
\[
\sum_{u\in V}\delta_u\,m_u=2n.
\]
\end{lemma}

\begin{proof}
The linear map $X\mapsto QX+XQ$ on symmetric matrices has eigenvalues $(1+\lambda_i)+(1+\lambda_j)>0$, so it is
invertible and~\eqref{eq:syl} determines $C$ from $D$. Since $e^{-tQ}\to0$ exponentially fast, the matrix
$\int_0^\infty e^{-tQ}De^{-tQ}\,dt$ is well defined, and it solves~\eqref{eq:syl}, because
$Q(\cdot)+(\cdot)Q$ applied to the integrand is $-\frac{d}{dt}\bigl(e^{-tQ}De^{-tQ}\bigr)$. Hence
\[
C=\int_0^\infty e^{-tQ}De^{-tQ}\,dt,\qquad
1=c(v,v)=\sum_u\delta_u\int_0^\infty\bigl(e^{-tQ}\bigr)_{vu}^2\,dt.
\]
Summing over $v$ and using the symmetry of $e^{-tQ}$,
\[
n=\sum_u\delta_u\int_0^\infty\bigl(e^{-2tQ}\bigr)_{uu}\,dt=\frac12\sum_u\delta_u\bigl(Q^{-1}\bigr)_{uu}.\qedhere
\]
\end{proof}

\section{The diagonal of \texorpdfstring{$(I+P)^{-1}$}{(I+P)\^{}(-1)}}

\begin{lemma}\label{lem:diag}
For every vertex $u$,
\[
\bigl(Q^{-1}\bigr)_{uu}\ge\frac{2n-3}{2(n-2)},
\]
with equality if and only if $u$ is adjacent to every other vertex.
\end{lemma}

\begin{proof}
Write $e_u=a+b$ with $a=\one/n$ and $b=e_u-\one/n\perp\one$. Since $\one$ is an eigenvector of $Q$ with eigenvalue $2$,
the cross terms vanish and
\[
\bigl(Q^{-1}\bigr)_{uu}=a^{\top}Q^{-1}a+b^{\top}Q^{-1}b=\frac1{2n}+b^{\top}Q^{-1}b .
\]
By the Cauchy--Schwarz inequality applied to $Q^{1/2}b$ and $Q^{-1/2}b$,
$(b^{\top}b)^2\le(b^{\top}Qb)(b^{\top}Q^{-1}b)$, with equality if and only if $Qb$ is a multiple of $b$. Here
$b^{\top}b=1-\frac1n$, and since $P(u,u)=0$,
\[
b^{\top}Qb=e_u^{\top}Qe_u-a^{\top}Qa=1-\frac2n>0 .
\]
Therefore
\[
\bigl(Q^{-1}\bigr)_{uu}\ge\frac1{2n}+\frac{(n-1)^2}{n(n-2)}=\frac{2n-3}{2(n-2)} .
\]
For the case of equality, $Pb=\frac1d\one_{N(u)}-\frac1n\one$, where $N(u)$ is the neighbourhood of $u$. If $Pb=\mu b$,
then the coordinate $u$ gives $-\frac1n=\mu(1-\frac1n)$, so $\mu=-\frac1{n-1}$, while a coordinate $w\notin
N(u)\cup\{u\}$ would give $-\frac1n=-\frac{\mu}n$, so $\mu=1$. Hence equality forces $N(u)=V\setminus\{u\}$.
Conversely, if $u$ is adjacent to all other vertices then $d=n-1$, $G=K_n$, and one checks directly that
$Pb=-\frac1{n-1}b$.
\end{proof}

Using only $(Q^{-1})_{uu}\ge1/Q_{uu}=1$ in place of Lemma~\ref{lem:diag}, the argument below gives $\Var S\le n/2$,
which is Proposition~14.22 of~\cite{AF} with a non-strict inequality. The gain comes from treating the constant vector,
on which $Q$ acts by $2$ rather than by something near $1$, separately.

\section{Proofs}

\begin{proof}[Proof of Theorem~\ref{thm:main}]
Put $m^*=\frac{2n-3}{2(n-2)}$. By Lemmas~\ref{lem:rep} and~\ref{lem:diag}, and $\delta_u\ge0$,
\[
2n=\sum_u\delta_um_u\ge m^*\sum_u\delta_u=4m^*\Var S ,
\]
so $\Var S\le n/(2m^*)=n(n-2)/(2n-3)$. If $G\ne K_n$, then $G$, being regular, has no vertex adjacent to all others, so
$m_u>m^*$ for every $u$ by Lemma~\ref{lem:diag}; and $\sum_u\delta_u>0$ by Lemma~\ref{lem:rep}. Hence the inequality is
strict. If $G=K_n$, every $m_u$ equals $m^*$ and equality holds, in agreement with~\eqref{eq:Kn}.
\end{proof}

\begin{proof}[Proof of Corollary~\ref{cor:vt}]
The linear system of Lemma~\ref{lem:eq} (the off-diagonal entries of~\eqref{eq:syl} together with $c(v,v)=1$) has
at most one solution. Indeed, the difference $E$ of two solutions has zero diagonal and $QE+EQ=F$ diagonal, so by the
proof of Lemma~\ref{lem:rep} $E=\int_0^\infty e^{-tQ}Fe^{-tQ}\,dt$, and its diagonal is $Mf$, where $f$ is the diagonal
of $F$ and $M=\int_0^\infty e^{-tQ}\circ e^{-tQ}\,dt$. By the Schur product theorem $M$ is positive definite, so $f=0$
and $E=0$. Every automorphism $\sigma$ of $G$ preserves the antivoter dynamics, so $c(\sigma v,\sigma w)$ also solves
the system. Hence $c$ is
automorphism-invariant, and on a vertex-transitive graph both $\delta_u=\delta$ and $m_u=\frac1n\tr Q^{-1}$ are
constant. Lemma~\ref{lem:rep} gives $\delta\tr Q^{-1}=2n$, and Lemma~\ref{lem:eq} gives
$\Var S=\frac n4\delta=\frac{n^2}{2\tr Q^{-1}}$, with $\tr Q^{-1}=\sum_i(1+\lambda_i)^{-1}$.
\end{proof}

\section{Numerical check}

The correlation equations of Lemma~\ref{lem:eq} form a linear system in the $\binom n2$ unknowns $c(v,w)$, $v\ne w$,
which we solved exactly for every connected non-bipartite regular graph on at most seven vertices, excluding cycles
(Table~\ref{tab:small}), and for $297$ random connected non-bipartite regular graphs with $5\le n\le14$ and $3\le d\le
n-1$. In every case $\Var S\le n(n-2)/(2n-3)$ with equality exactly for $K_n$, and every diagonal entry of
$(I+P)^{-1}$ was at least $m^*$, strictly so when $d<n-1$. The formula of Corollary~\ref{cor:vt} agrees with the exact
solution to within $4\cdot10^{-15}$ on eight non-bipartite vertex-transitive graphs, among them the Petersen graph
($\Var S=40/13$), $C_3\times C_3$, two circulants and the cycle $C_7$ (where $\Var S=1$).

\begin{table}[ht]
\centering
\begin{tabular}{rrrr}
\toprule
$n$ & graph & $d$ & $\Var S$\\
\midrule
4 & $K_4$ & 3 & 1.6000\\
5 & $K_5$ & 4 & 2.1429\\
6 & $K_6$ & 5 & 2.6667\\
6 & octahedron & 4 & 2.4000\\
6 & prism $C_3\times K_2$ & 3 & 1.9459\\
7 & $K_7$ & 6 & 3.1818\\
7 & $4$-regular (two graphs) & 4 & 2.7067,\ 2.4653\\
\bottomrule
\end{tabular}
\caption{Exact stationary variance of $S$ on every connected non-bipartite regular graph with at most seven vertices,
cycles excluded.}\label{tab:small}
\end{table}

\section*{Acknowledgements and funding}
The research presented in this paper was supported by the European Research
Council (ERC) under the European Union's Horizon Europe research and
innovation programme (grant agreement No.~101041711), by the Simons
Foundation as part of the Collaboration on the Mathematical and Scientific
Foundations of Deep Learning, by Heights Labs, by Convex Nexus Capital, by
the Israel Science Foundation (grant number 2258/19), and by the Israel
Science Foundation (ISF Grant 4101/25). Computational resources were
provided in part by Google Cloud through its education credits program
(credit 530550103).

This work was done by multiple AI agents, Claude (Anthropic) and Codex (OpenAI), connected to Math Brain~\cite{MathBrainRepo} and working together inside the Brain framework~\cite{BrainFramework}, a shared, persistent memory in which the agents record claims, proofs, failed attempts and methods, retrieve one another's results, coordinate who works on what, and audit each other's work. For more on the details behind the Brain and Math Brain, see the Math Brain paper~\cite{MathBrainPaper}, the Math Brain repository~\cite{MathBrainRepo} and the Brain framework repository~\cite{BrainFramework}.

\section*{Declaration of competing interest}
The author holds equity in Heights Labs and in Convex Nexus Capital and
has a management role at Convex Nexus Capital. These entities are included
in the funding declaration above.

\section*{Declaration of generative AI and AI-assisted technologies}
The paper comes out of Math Brain, an ongoing project on the use of AI systems for research mathematics. In November 2026
Claude Opus~5.5 (Anthropic), working in Claude Code, found the proofs, wrote and ran the computations of
Table~\ref{tab:small}, and wrote the paper. The author directed the work and is responsible for the content of the
paper and for decisions concerning its dissemination.

\section*{Data availability}
No data were used. The code for Table~\ref{tab:small} and for the random-graph check, with its output, is available at
\url{https://github.com/yspennstate/antivoter-variance-complete-graph}.

The paper is part of the Math Brain project, whose public repository is \url{https://github.com/yspennstate/math-brain-public}~\cite{MathBrainRepo}.

\begin{thebibliography}{9}
\bibitem{AF} D.~Aldous and J.~A.~Fill, \emph{Reversible Markov Chains and Random Walks on Graphs}, unfinished
monograph, 2002 (recompiled 2014), \url{https://www.stat.berkeley.edu/~aldous/RWG/book.html}, Section~14.4.
\bibitem{DW} P.~Donnelly and D.~Welsh, \emph{The antivoter problem: random $2$-colourings of graphs}, in Graph Theory
and Combinatorics (Cambridge, 1983), Academic Press, London, 1984.
\bibitem{RR} Y.~Rinott and V.~Rotar, \emph{On coupling constructions and rates in the CLT for dependent summands with
applications to the antivoter model and weighted $U$-statistics}, Ann. Appl. Probab. \textbf{7} (1997), 1080--1105.
\bibitem{Ross} N.~Ross, \emph{Fundamentals of Stein's method}, Probab. Surv. \textbf{8} (2011), 210--293.
\bibitem{MathBrainRepo} Y.~Shmalo, \emph{Math Brain: public snapshot, code and tests}, 2026, \url{https://github.com/yspennstate/math-brain-public}.

\bibitem{MathBrainPaper} Y.~Shmalo, \emph{Math Brain: shared, auditable memory for multi-agent mathematical research}, October 2026, \url{https://github.com/yspennstate/math-brain-public/tree/main/papers/math-brain-2026-10-24}.

\bibitem{BrainFramework} Y.~Shmalo, \emph{Brain Framework}, public software repository, 2026, \url{https://github.com/yspennstate/brain-framework-public}.

\end{thebibliography}

\end{document}
